
\section{Conclusion}

This thesis presented a comprehensive study on the application of Deep Reinforcement Learning for the energy management of residential microgrids. By formulating the battery control problem as a Markov Decision Process and solving it using the Soft Actor-Critic (SAC) algorithm, this research demonstrated that data-driven agents can achieve near-optimal economic performance without requiring explicit models of the system dynamics or stochastic parameters.

The experimental results rigorously confirmed that the SAC agent significantly outperforms distinct baselines, including traditional rule-based heuristics and the on-policy Proximal Policy Optimization (PPO) algorithm. Specifically, the proposed SAC agent achieved a mean daily profit within 3.53\% of the theoretical upper bound established by a Linear Programming benchmark with perfect foresight. This narrow optimality gap indicates that the Maximum Entropy framework allows the agent to effectively explore the continuous action space and discover sophisticated arbitrage strategies that capitalize on Time-of-Use price differentials. Translating these daily economic gains into a long-term perspective, the SAC agent is estimated to yield annual savings of approximately \$7,128 compared to a standard rule-based heuristic controller. This substantial potential for cost reduction highlights the practical financial viability of deploying intelligent learning-based agents for residential energy management.

Furthermore, this study addressed the critical challenge of forecast uncertainty through comprehensive robustness analysis. When subjected to severe ``Sim-to-Real'' gaps characterized by 30\% correlated forecast errors, the PPO agent demonstrated exceptional robustness with a score of 90.18, retaining over 90\% of its original performance. The SAC and R-SAC agents achieved robustness scores of 80.14 and 80.21 respectively, substantially outperforming the LP benchmark which scored only 76.90. Importantly, despite its lower robustness, SAC maintained superior absolute profits under uncertainty due to its strong baseline performance.

These findings advocate for the adoption of entropy-regularized Reinforcement Learning as a robust, scalable, and economically viable solution for the next generation of smart grid energy management systems. The choice between SAC and PPO depends on operational priorities: SAC is recommended for environments with reasonably reliable forecasts where profit maximization is paramount, while PPO may be preferred in highly uncertain conditions where operational stability is the primary concern.

